\label{chap:benchmark}

Este capítulo presenta la metodología experimental utilizada para evaluar y comparar
las diferentes alternativas criptográficas consideradas en este trabajo. El objetivo
no consiste únicamente en determinar cuál implementación presenta el menor tiempo de
ejecución, sino en estudiar los compromisos existentes entre rendimiento, consumo de
recursos, comunicación, seguridad, complejidad de implementación y requisitos
operacionales.

% ============================================================
\subsection{Objetivos de la evaluación}
\label{sec:benchmark-objectives}

La evaluación experimental persigue los siguientes objetivos:

\begin{enumerate}
    \item Comparar el costo computacional de los diferentes protocolos durante las
    principales etapas de autenticación.

    \item Comparar el consumo de memoria y almacenamiento requerido por cada
    alternativa.

    \item Determinar el volumen de información que debe transmitirse durante una
    autenticación y estudiar su comportamiento ante diferentes condiciones de red.

    \item Analizar las condiciones de seguridad asociadas con los parámetros utilizados
    por cada protocolo.

    \item Examinar los principales ataques relevantes para cada familia criptográfica
    y las condiciones necesarias para mitigarlos.

    \item Comparar la complejidad de implementación, configuración y despliegue de cada
    alternativa.

    \item Analizar las dependencias de confianza y los requisitos de infraestructura de
    cada protocolo.

    \item Evaluar el impacto que tendría cada alternativa sobre la experiencia de un
    usuario final.

    \item Identificar los principales compromisos entre seguridad, rendimiento,
    comunicación, memoria y facilidad de despliegue.

    \item Comparar el protocolo propuesto en este trabajo con las demás alternativas
    bajo las mismas condiciones experimentales.
\end{enumerate}


% ============================================================
\subsection{Principio de comparabilidad}
\label{sec:benchmark-comparability}

Para garantizar una comparación significativa, todos los candidatos implementan una
misma funcionalidad abstracta de autenticación.

El flujo experimental se divide conceptualmente en dos momentos: registro
(\textit{enrollment}) y autenticación.

\subsubsection{Registro}

Durante el registro se ejecutan las operaciones necesarias para preparar al usuario y
al verificador antes de realizar una autenticación.

De manera abstracta, esta etapa se representa como:

\begin{equation}
    (pk, sk, \theta) \leftarrow \operatorname{Setup}()
\end{equation}

donde:

\begin{itemize}
    \item $pk$ representa la información pública necesaria para verificar al usuario;
    \item $sk$ representa el material secreto conservado por el usuario o demostrador;
    \item $\theta$ representa parámetros adicionales requeridos por el protocolo.
\end{itemize}

Dependiendo de la construcción, esta operación puede incluir generación de claves,
muestreo de vectores, construcción de matrices, generación de parámetros públicos o
precomputaciones.

\subsubsection{Autenticación}

Una autenticación se modela mediante las siguientes operaciones:

\begin{align}
    c &\leftarrow \operatorname{Challenge}(pk), \\
    r &\leftarrow \operatorname{Respond}(sk,c), \\
    b &\leftarrow \operatorname{Verify}(pk,c,r),
\end{align}

donde:

\begin{itemize}
    \item $c$ es el reto generado por el verificador;
    \item $r$ es la respuesta generada por el usuario;
    \item $b \in \{0,1\}$ representa el resultado de la verificación.
\end{itemize}

Para una ejecución válida se espera que:

\begin{equation}
    \operatorname{Verify}(pk,c,
    \operatorname{Respond}(sk,c)) = 1.
\end{equation}

Esta abstracción permite medir las diferentes construcciones utilizando las mismas
etapas funcionales aun cuando las operaciones matemáticas internas sean diferentes.

\subsubsection{Separación entre registro y autenticación}

El tiempo de generación de claves no se incorpora automáticamente al tiempo de cada
autenticación, puesto que normalmente corresponde a una operación realizada durante
el registro inicial.

Por lo tanto, se reportan por separado:

\begin{equation}
    T_{\text{enrollment}}
\end{equation}

y

\begin{equation}
    T_{\text{authentication}}
    =
    T_{\text{challenge}}
    +
    T_{\text{response}}
    +
    T_{\text{verification}}.
\end{equation}

Cuando una implementación requiera operaciones adicionales en alguna de estas etapas,
estas deberán documentarse y contabilizarse explícitamente.


% ============================================================
\subsection{Protocolos evaluados}
\label{sec:benchmark-protocols}

La Tabla~\ref{tab:protocol-summary} resume las alternativas estudiadas.

\begin{table}[H]
\centering
\caption{Protocolos incluidos en la evaluación experimental.}
\label{tab:protocol-summary}
\small
\begin{tabularx}{\textwidth}{lYYY}
\toprule
\textbf{Protocolo} &
\textbf{Fundamento matemático} &
\textbf{Tipo de construcción} &
\textbf{Rol en la comparación} \\
\midrule

ECDSA &
Problema del logaritmo discreto sobre curvas elípticas &
Criptografía tradicional &
Referencia de industria \\

LWE estándar &
Learning With Errors &
Retículos &
Candidato post-cuántico \\

Binary-LWE &
Learning With Errors con secreto restringido &
Retículos &
Variante LWE \\

Ring-LWE &
Learning With Errors sobre estructuras algebraicas de anillos &
Retículos estructurados &
Variante estructurada \\

LWR &
Learning With Rounding &
Retículos &
Variante basada en redondeo \\

LWE propuesto &
\todoresult{describir fundamento exacto} &
Retículos &
Propuesta de este trabajo \\

\bottomrule
\end{tabularx}
\end{table}

ECDSA se utiliza como referencia de una alternativa ampliamente adoptada en sistemas
criptográficos tradicionales. Las construcciones basadas en LWE y problemas
relacionados permiten estudiar alternativas fundamentadas en retículos y analizar sus
propiedades frente a los requisitos considerados en este proyecto.

La comparación se realiza utilizando un mismo modelo funcional de autenticación y un
entorno experimental controlado. De esta manera, las diferencias reportadas buscan
reflejar las características de cada alternativa criptográfica y no diferencias
artificiales introducidas por procedimientos de medición distintos.